% TOP_PROBLEM: {"id": 7200083, "title": "Is the upper bound of a closed knot's minimum ropelength linear to its crossing number", "category_id": 6, "category": {"id": 6, "name": "geometry", "display_name": "Geometry", "description": "Euclidean and non-Euclidean geometry, geometric structures.", "slug": "geometry", "order_index": 6, "created_at": "2026-07-31T15:26:25.671Z"}, "status": "partially_solved"}
% FIRST_POSTED: null
% ENTRY_KIND: statement_only
% Imported from: batch20.tex; UnsolvedMath ID 7200083
\documentclass[11pt]{article}
\usepackage[a4paper,margin=25mm]{geometry}
\usepackage{amsmath,amssymb,mathtools}
\usepackage{fontspec,unicode-math}
\setmainfont{Latin Modern Roman}
\setsansfont{Latin Modern Sans}
\setmathfont{Latin Modern Math}
\usepackage{xurl,hyperref}
\hypersetup{colorlinks=true,linkcolor=blue,urlcolor=blue,pdftitle={UnsolvedMath Problem 7200083}}
\newcommand{\sref}[2]{\hyperlink{source:#1:#2}{[S#2]}}
\newcommand{\cataloguescope}[1]{\paragraph{Mathematical question.}#1}
\begin{document}
% UnsolvedMath ID 7200083; source catalogue code AMR-071-0083
\setcounter{section}{7200082}
\section{Is the upper bound of a closed knot's minimum ropelength linear to its crossing number}\label{7200083}
{\small\sffamily UnsolvedMath ID 7200083 · Geometry}
\subsection{Definitions and mathematical statement}

\paragraph{Shared notation.}$\mathbb N=\{1,2,\ldots\}$, and $\mathbb Z,\mathbb Q,\mathbb R,\mathbb C$ denote the integers, rationals, reals and complex numbers. Any other conventions are specified locally.


A knot type $K$ is an ambient-isotopy class of embedded closed curves in $\mathbb R^3$. For a $C^{1,1}$ representative $\gamma$ (continuously differentiable with Lipschitz derivative), its thickness $\tau(\gamma)$ is the supremal radius of disjoint normal disks forming an embedded tube. Its ropelength is $\operatorname{length}(\gamma)/\tau(\gamma)$, and $L(K)$ is the infimum over all representatives. This is the radius convention for unit thickness. The crossing number $\operatorname{Cr}(K)$ is the minimum number of transverse crossings in a regular planar knot diagram.
\paragraph{Statement recorded in the supplied source.}
Does there exist a constant $C>0$, independent of the knot type, such that $L(K)\leq C\max\{1,\operatorname{Cr}(K)\}$ for every knot type $K$? Equivalently, is $L(K)=O(\operatorname{Cr}(K))$ uniformly over nontrivial knot types? \sref{7200083}{2}
\subsection{Short English statement}
Can every nontrivial knot be tied with unit-radius rope whose length is bounded by a universal constant times its crossing number?
\subsection{Sources}
\begin{description}
\item[\hypertarget{source:7200083:1}{S1}] ULAM.AI and the UnsolvedMath Contributors, UnsolvedMath: A Curated Collection of Open Mathematics Problems, record 7200083 (AMR-071-0083). CC BY 4.0. \url{https://huggingface.co/datasets/ulamai/UnsolvedMath}.
\item[\hypertarget{source:7200083:2}{S2}] Yuanan Diao, Claus Ernst, Attila Por and Uta Ziegler, The Ropelengths of Knots Are Almost Linear in Terms of Their Crossing Numbers (2009), v1; arXiv:0912.3282. \url{https://arxiv.org/pdf/0912.3282}.
\end{description}
\label{end:7200083}
\end{document}
